\documentclass[a4paper,11pt,fleqn]{scrartcl}
\usepackage{../../Styles/Eval}

\newcommand{\chnum}{Chapitre 6 et 21}
\newcommand{\chtitle}{Transformations chimiques et électricité}
\newcommand{\dtype}{Corrigé DS}

\begin{document}
	\maketitle
	
	\section{Autour des définitions}
    
    Equation de réaction :
    
    \[
    CH_4(g) + 2\,O_2(g) \rightarrow CO_2(g) + 2\,H_2O(g)
    \]
    
    \begin{enumerate}
        \item Que désignent chacun des éléments encadrés ? (3pts)
        
        \begin{itemize}
            \item $O_2(g)$ : le dioxygène, réactif à l’état gazeux.
            \item $\rightarrow$ : signifie « se transforme en » ou « donne ».
            \item $CO_2$ : le dioxyde de carbone.
            \item $(g)$ : indique l’état gazeux.
            \item Le coefficient 2 devant $H_2O$ : il y a deux molécules d’eau produites.
        \end{itemize}

        \item Traduire cette équation de réaction en langage courant. (2pts)
        
        Le méthane réagit avec le dioxygène pour former du dioxyde de carbone et de l’eau.
        
    \end{enumerate}

    \section{Équilibrer des équations chimiques}

    \begin{enumerate}
        \item
        
        \[
        C_6H_{12}O_6(s) + 6\,O_2(g) \rightarrow 6\,CO_2(g) + 6\,H_2O(l)
        \]

        \item
        
        \[
        4\,Fe(s) + 3\,O_2(g) \rightarrow 2\,Fe_2O_3(s)
        \]

        \item
        
        \[
        N_2H_4(l) + O_2(g) \rightarrow N_2(g) + 2\,H_2O(l)
        \]

        \item
        
        \[
        2\,Al^{3+}(aq) + 3\,Cu(s) \rightarrow 2\,Al(s) + 3\,Cu^{2+}(aq)
        \]
        
    \end{enumerate}

    \section{Lois des circuits électriques}

    \begin{enumerate}
        \item Représenter sur le schéma les tensions orientées $U_L$, $U_{R1}$ et $U_{R2}$. (1pt)
        
        Les flèches des tensions sont orientées dans le sens opposé du courant dans les dipôles récepteurs :
        
        \begin{itemize}
            \item $U_L$ : orientée de la borne de sortie du courant vers la borne d’entrée de la lampe.
            \item $U_{R1}$ : orientée dans le sens opposé à $I_2$.
            \item $U_{R2}$ : orientée dans le sens opposé à $I_3$.
        \end{itemize}

        \item Calculer la valeur de $I_2$. (1pt)
        
        Loi des nœuds :
        
        \[
        I_1 = I_2 + I_3
        \]
        
        Donc :
        
        \[
        I_2 = I_1 - I_3
        \]
        
        \[
        I_2 = 0{,}10 - 0{,}04
        \]
        
        \[
        I_2 = 0{,}06~A
        \]

        \item Calculer la valeur de $R_1$. (1pt)
        
        La tension aux bornes de $R_1$ est égale à celle du générateur :
        
        \[
        U_{R1} = 12~V
        \]
        
        Loi d’Ohm :
        
        \[
        R_1 = \frac{U_{R1}}{I_2}
        \]
        
        \[
        R_1 = \frac{12}{0{,}06}
        \]
        
        \[
        R_1 = 200~\Omega
        \]
        
    \end{enumerate}

    \section{Quantité de matière}

Dans la réaction de combustion du méthanol :

\[
CH_3OH(l) + O_2(g) \rightarrow CO_2(g) + H_2O(l)
\]

on introduit $20g$ de méthanol et $20g$ de dioxygène.

\begin{enumerate}
    \item Equilibrer l’équation de réaction. (1pt)

    \[
    2\,CH_3OH(l) + 3\,O_2(g) \rightarrow 2\,CO_2(g) + 4\,H_2O(l)
    \]

    \item Calculer la quantité de matière de méthanol et de dioxygène introduits. (2pt)

    \underline{Méthanol}

    Masse d’une molécule de méthanol :

    \[
    m_{CH_3OH}=1\times m_C+4\times m_H+1\times m_O
    \]

    \[
    m_{CH_3OH}=1,99\times10^{-26}+4\times1,67\times10^{-27}+2,66\times10^{-26}
    \]

    \[
    m_{CH_3OH}=5,318\times10^{-26}~kg
    \]

    Conversion de la masse introduite :

    \[
    20g = 2,0\times10^{-2}~kg
    \]

    Nombre de molécules dans l’échantillon :

    \[
    N=\frac{m}{m_{\text{molécule}}}
    \]

    \[
    N(CH_3OH)=\frac{2,0\times10^{-2}}{5,318\times10^{-26}}
    \]

    \[
    N(CH_3OH)=3,76\times10^{23}
    \]

    Quantité de matière :

    \[
    n=\frac{N}{N_A}
    \]

    \[
    n(CH_3OH)=\frac{3,76\times10^{23}}{6,02\times10^{23}}
    \]

    \[
    n(CH_3OH)=0,625~mol
    \]

    \underline{Dioxygène}

    Masse d’une molécule de dioxygène :

    \[
    m_{O_2}=2\times m_O
    \]

    \[
    m_{O_2}=2\times2,66\times10^{-26}
    \]

    \[
    m_{O_2}=5,32\times10^{-26}~kg
    \]

    Nombre de molécules dans l’échantillon :

    \[
    N(O_2)=\frac{2,0\times10^{-2}}{5,32\times10^{-26}}
    \]

    \[
    N(O_2)=3,76\times10^{23}
    \]

    Quantité de matière :

    \[
    n(O_2)=\frac{3,76\times10^{23}}{6,02\times10^{23}}
    \]

    \[
    n(O_2)=0,625~mol
    \]

    \item Le premier réactif à être consommé entièrement est nommé \underline{réactif limitant}. Déterminer le réactif limitant de cette réaction. (1pt)

    D’après l’équation :

    \[
    2\,CH_3OH + 3\,O_2
    \]

    il faut $1,5$ fois plus de dioxygène que de méthanol.

    Pour $0,625~mol$ de méthanol, il faudrait :

    \[
    n(O_2)=0,625\times1,5
    \]

    \[
    n(O_2)=0,938~mol
    \]

    Or on ne dispose que de $0,625~mol$ de dioxygène.

    Le dioxygène est donc le réactif limitant.

    \item En déduire la quantité de matière de dioxyde de carbone produite. (2pt)

    D’après l’équation :

    \[
    3\,O_2 \rightarrow 2\,CO_2
    \]

    Donc :

    \[
    n(CO_2)=\frac{2}{3}\times n(O_2)
    \]

    \[
    n(CO_2)=\frac{2}{3}\times0,625
    \]

    \[
    n(CO_2)=0,417~mol
    \]

    \item Calculer la masse de dioxyde de carbone produite. (2pt)

    Masse d’une molécule de dioxyde de carbone :

    \[
    m_{CO_2}=1\times m_C+2\times m_O
    \]

    \[
    m_{CO_2}=1,99\times10^{-26}+2\times2,66\times10^{-26}
    \]

    \[
    m_{CO_2}=7,31\times10^{-26}~kg
    \]

    Nombre de molécules produites :

    \[
    N(CO_2)=n\times N_A
    \]

    \[
    N(CO_2)=0,417\times6,02\times10^{23}
    \]

    \[
    N(CO_2)=2,51\times10^{23}
    \]

    Masse produite :

    \[
    m=N\times m_{\text{molécule}}
    \]

    \[
    m(CO_2)=2,51\times10^{23}\times7,31\times10^{-26}
    \]

    \[
    m(CO_2)=1,84\times10^{-2}~kg
    \]

    \[
    m(CO_2)\approx18,4~g
    \]

\end{enumerate}
\end{document}